% ===============================================================
% CONCLUSION GÉNÉRALE
% ===============================================================
\vspace*{-3cm}
\chapter*{Conclusion générale}
\addcontentsline{toc}{chapter}{Conclusion générale}
\vspace*{-1cm}

Ce projet de fin d'année nous a permis de concevoir et développer WC Fantasy, une application mobile de fantasy football dédiée à la Coupe du Monde 2026. Partant d'un vide identifié dans l'offre existante --- l'absence d'une plateforme mobile combinant fantasy sport, fonctionnalités sociales en temps réel et gamification pour cet événement mondial --- nous avons proposé une solution complète et opérationnelle.

Sur le plan technique, le projet a mobilisé un large spectre de compétences : architecture backend avec Node.js/Express/PostgreSQL, développement mobile cross-platform avec Flutter/Dart, communications temps réel avec Socket.io, gestion du cache avec Redis, intégration d'APIs tierces (football-data.org, Google Gemini, Brevo) et sécurisation multicouche (JWT, bcrypt, rate limiting, transactions ACID). L'ensemble représente 24 routes API REST, 20 tables en base de données, et plusieurs modules spécialisés (calcul de points, synchronisation de scores, email, chatbot).

Les compétences acquises au cours de ce projet sont multiples : maîtrise de l'architecture client-serveur découplée, gestion des états asynchrones complexes, conception de schémas de base de données relationnels, sécurisation d'une API REST, développement mobile réactif avec Riverpod, et gestion de projet en autonomie avec respect d'un cahier des charges structuré.



\textbf{Perspectives d'amélioration :} plusieurs évolutions sont envisagées pour les versions futures. Premièrement, le déploiement en production sur une infrastructure cloud (AWS ou Railway) avec containerisation Docker et reverse proxy Nginx. Deuxièmement, l'activation du microservice Python de recommandations ML de joueurs basées sur les statistiques historiques. Troisièmement, l'intégration complète des notifications push Firebase (FCM pour Android, APNs pour iOS). Quatrièmement, l'ajout d'un mode spectateur permettant de suivre des équipes d'autres utilisateurs sans en être membre. Enfin, la création d'une interface web (Next.js) complémentaire à l'application mobile pour couvrir l'ensemble des plateformes.

WC Fantasy représente une démonstration concrète de notre capacité à concevoir et réaliser un système logiciel complet, de la spécification à la mise en \oe{}uvre, en appliquant les meilleures pratiques du développement web et mobile modernes.
